%%FILE blueprint/src/macros/common.tex
% In this file you should put all LaTeX macros to be used
% both by the pdf version and the web version.
% This should be most of your macros.

% Put any macro needed for the project here. Those macros for both the web and print versions of the blueprint.

\newcommand{\Z}{\mathbb{Z}}
\newcommand{\bZ}{\Z}
\newcommand{\N}{\mathbb{N}}
\newcommand{\bN}{\N}
\newcommand{\Q}{\mathbb{Q}}
\newcommand{\bQ}{\Q}
\newcommand{\R}{\mathbb{R}}
\newcommand{\bR}{\R}
\newcommand{\C}{\mathbb{C}}
\newcommand{\bC}{\C}
\newcommand{\Zpos}{\Z_{>0}}
\newcommand{\Znn}{\Z_{\ge 0}}
\newcommand{\Npos}{\N_{+}}

\newcommand{\PR}{\mathsf{P}}
\newcommand{\EX}{\mathsf{E}}

\newcommand{\indof}[1]{\mathbb{I}_{{#1}}}
\newcommand{\PRmeas}{\mu}

\newcommand{\eps}{\varepsilon}

\newcommand{\oo}{\mathfrak{o}}
\newcommand{\OO}{\mathcal{O}}

\DeclareMathOperator{\Aut}{Aut}
\DeclareMathOperator{\Hom}{Hom}
\DeclareMathOperator{\GL}{GL}
\DeclareMathOperator{\PGL}{PGL}
\DeclareMathOperator{\PSL}{PSL}
\DeclareMathOperator{\SL}{SL}
\DeclareMathOperator{\tr}{tr}
\DeclareMathOperator{\id}{id}
\DeclareMathOperator{\dmn}{dim}

\newcommand{\set}[1]{\left\{ #1 \right\}}
%\newcommand{\idOf}[1]{\id_{#1}}

%\newcommand{\lcinv}[1]{\vec{#1}^{-1}}
%\newcommand{\rcinv}[1]{\cev{#1}^{-1}}
\newcommand{\lcinv}[1]{{#1}^{\rightarrow 1}}
\newcommand{\rcinv}[1]{{#1}^{\leftarrow 1}}

\newcommand{\AffOf}[1]{\mathrm{Aff}_{#1}}
\newcommand{\AffR}{\AffOf{\R}}
\newcommand{\OriAffR}{\mathrm{Aff}^+_{\R}}

\newcommand{\tendsinlaw}{\overset{\mathrm{d}}{\longrightarrow}}

\newcommand{\Gumbel}{\mathrm{Gumbel}}
\newcommand{\Weibull}{\mathrm{Weibull}}
\newcommand{\Frechet}{\mathrm{Fréchet}}
\newcommand{\GumbelCDF}{\Lambda}
\newcommand{\WeibullCDF}[1]{\Psi_{#1}}
\newcommand{\FrechetCDF}[1]{\Phi_{#1}}
%\newcommand{\GumbelCDF}{F_{\Gumbel}}
%\newcommand{\WeibullCDF}[1]{F^{(#1)}_{\Weibull}}
%\newcommand{\FrechetCDF}[1]{F^{(#1)}_{\Frechet}}

\newcommand{\extendcdf}[1]{\widetilde{#1}}
\newcommand{\oneDivOneSub}[1]{\frac{\mathbf{1}}{\mathbf{1}-\extendcdf{#1}}}
\newcommand{\oneDivNegLog}[1]{\frac{\mathbf{1}}{\widetilde{\log} \big( 1 / \extendcdf{#1} \big)}}

\newcommand{\ud}{\mathrm{d}}
\newcommand{\der}[1]{\frac{\ud}{\ud #1}}
\newcommand{\dder}[1]{\frac{\ud^2}{\ud {#1}^2}}

\newcommand{\Leb}{\Lambda}%{\mu_{\mathrm{Leb}}}

%%FILE blueprint/src/macros/web.tex
% In this file you should put macros to be used only by
% the web version. Of course they should have a corresponding
% version in macros/print.tex.
% Typically the printed version could have more fancy decorations.
% This will probably be a very short file.

\newtheorem{theorem}{Theorem}[chapter]
\newtheorem{proposition}[theorem]{Proposition}
\newtheorem{lemma}[theorem]{Lemma}
\newtheorem{sublemma}[theorem]{Sub-lemma}
\newtheorem{corollary}[theorem]{Corollary}

\theoremstyle{definition}
\newtheorem{definition}[theorem]{Definition}

\theoremstyle{remark}
\newtheorem{remark}[theorem]{Remark}
\newtheorem{example}[theorem]{Example}

%%FILE blueprint/src/macros/print.tex
% In this file you should put macros to be used only by
% the printed version. Of course they should have a corresponding
% version in macros/web.tex.
% Typically the printed version could have more fancy decorations.
% This should be a very short file.
%
% This file starts with dummy macros that ensure the pdf
% compiler will ignore macros provided by plasTeX that make
% sense only for the web version, such as dependency graph
% macros.

%\declaretheorem[numberwithin=chapter]{theorem}
%\declaretheorem[sibling=theorem]{proposition}
%\declaretheorem[sibling=theorem]{corollary}
%\declaretheorem[sibling=theorem]{remark}
%\declaretheorem[sibling=theorem]{lemma}
%\declaretheorem[sibling=theorem]{definition}
%\declaretheorem[sibling=theorem]{example}

\newtheorem{theorem}{Theorem}[chapter]
\newtheorem{proposition}[theorem]{Proposition}
\newtheorem{lemma}[theorem]{Lemma}
\newtheorem{sublemma}[theorem]{Sub-lemma}
\newtheorem{corollary}[theorem]{Corollary}

\theoremstyle{definition}
\newtheorem{definition}[theorem]{Definition}

\theoremstyle{remark}
\newtheorem{remark}[theorem]{Remark}
\newtheorem{example}[theorem]{Example}

% We neutralise the Plastex commands
\newcommand{\uses}[1]{}
\newcommand{\proves}[1]{}
\newcommand{\discussion}[1]{}
\newcommand{\lean}[1]{}
\newcommand{\leanok}{}
\newcommand{\mathlibok}{}
\newcommand{\notready}{}

% % Dummy macros that make sense only for web version.
% \newcommand{\lean}[1]{}
% \newcommand{\discussion}[1]{}
% \newcommand{\leanok}{}
% \newcommand{\mathlibok}{}
% \newcommand{\notready}{}
% % Make sure that arguments of \uses and \proves are real labels, by using invisible refs:
% % latex prints a warning if the label is not defined, but nothing is shown in the pdf file.
% % It uses LaTeX3 programming, this is why we use the expl3 package.
% \ExplSyntaxOn
% \NewDocumentCommand{\uses}{m}
%  {\clist_map_inline:nn{#1}{\vphantom{\ref{##1}}}%
%   \ignorespaces}
% \NewDocumentCommand{\proves}{m}
%  {\clist_map_inline:nn{#1}{\vphantom{\ref{##1}}}%
%   \ignorespaces}
% \ExplSyntaxOff
